\subsection{Bifurcation Theory}
% Saddle-node, pitchfork, Hopf bifurcation, normal forms.

A parameter can change not merely the position of equilibria but also their number and stability. The harvested logistic equation is a complete scalar model of such a transition.

\begin{definition} \label{def:scalar-saddle-node-pattern} \hlt{Scalar Saddle-Node Pattern}\\
Let $\Lambda\subseteq\R$ be an interval, $F:\R\times\Lambda\to\R$, and $F_h(x)\defeq F(x,h)$.\\
For $h_c\in\Lambda$ and $x_c\in\R$, say $h_c$ exhibits a \hlt{saddle-node pattern} at $x_c$ iff $\exists\varepsilon>0$, a neighborhood $X$ of $x_c$,\\
and $\sigma\in\{-1,1\}$ with nonempty $H_\pm\defeq\{h\in\Lambda:0<\pm\sigma(h-h_c)<\varepsilon\}$ such that:
\begin{enumerate}[label=(\roman*)]
\setlength{\itemsep}{0pt}
\item $F_{h_c}^{-1}(0)\cap X=\{x_c\}$.
\item $\exists x_\pm:H_+\to X$ with $x_-(h)<x_+(h)$ and $F_h(x_\pm(h))=0$ $\forall h\in H_+$, and $\lim\limits_{h\to h_c,\ h\in H_+}x_\pm(h)=x_c$.
\item $F_h^{-1}(0)\cap X=\emptyset$ $\forall h\in H_-$.
\end{enumerate}
\hlt{Key Insight}: A saddle-node is local creation or annihilation of two equilibria.
\end{definition}

\begin{proposition} \label{prop:harvested-logistic-equilibria} \hlt{Equilibrium Branches of the Harvested Logistic Equation}\\
For $h\in[0,+\infty)$, let $F_h(x)\defeq(1-x)x-h$, $h_c\defeq\frac14$, and $x_c\defeq\frac12$.
\begin{enumerate}[label=(\roman*)]
\setlength{\itemsep}{0pt}
\item If $0\leq h<h_c$, then $F_h^{-1}(0)=\{x_-(h),x_+(h)\}$, where $x_\pm(h)\defeq\frac12(1\pm\sqrt{1-4h})$ and $x_-(h)<x_+(h)$.
\item If $h=h_c$, then $F_{h_c}^{-1}(0)=\{x_c\}$ and $F_{h_c}(x)=-(x-x_c)^2$.
\item If $h>h_c$, then $F_h^{-1}(0)=\emptyset$.
\item The branches $x_\pm:(0,h_c)\to\R$ coalesce at $x_c$: $x_\pm(h)\to x_c$ as $h\uparrow h_c$.
\end{enumerate}
\begin{proof}
For all $h\geq0$, $F_h(x)=h_c-h-(x-x_c)^2$, hence $F_h(x)=0$ iff $(x-x_c)^2=h_c-h$.
\begin{enumerate}[label=(\roman*)]
\setlength{\itemsep}{0pt}
\item If $0\leq h<h_c$, then $h_c-h>0$, so $F_h^{-1}(0)=\{x_c\pm\sqrt{h_c-h}\}=\{x_-(h),x_+(h)\}$.
\item If $h=h_c$, then $F_{h_c}(x)=-(x-x_c)^2$, so $F_{h_c}^{-1}(0)=\{x_c\}$ and the zero is double.
\item If $h>h_c$, then $h_c-h<0$, so $(x-x_c)^2=h_c-h$ has no real solution.
\item Since $x_\pm(h)=x_c\pm\sqrt{h_c-h}$ and $h_c-h\downarrow0$ as $h\uparrow h_c$, $x_\pm(h)\to x_c$.
\end{enumerate}
\hlt{Structural Commentary}: Completing the square isolates the zero count.\\
\hlt{Key Insight}: The critical value is exactly where the two roots merge.
\end{proof}
\end{proposition}

\begin{lemma} \label{lem:scalar-phase-line-barrier-limit} \hlt{Scalar Phase-Line Barriers and Limits}\\
Let $F\in C^1(\R,\R)$, and let $x:[t_0,T_+)\to\R$ be a maximal forward solution of $x'=F(x)$.
\begin{enumerate}[label=(\roman*)]
\setlength{\itemsep}{0pt}
\item If $e\in F^{-1}(0)$ and $x(\tau)=e$ for some $\tau\in[t_0,T_+)$, then $x\equiv e$ near $\tau$; no nonconstant branch crosses $e$.
\item If $x$ is monotone and bounded, then $\ell\defeq\lim\limits_{t\uparrow T_+}x(t)$ exists, $T_+=+\infty$, and $\ell\in F^{-1}(0)$.
\end{enumerate}
\begin{proof}
\begin{enumerate}[label=(\roman*)]
\setlength{\itemsep}{0pt}
\item Let $u\equiv e$. Then $u'=F(u)=0$ and $u(\tau)=x(\tau)$.\\
By uniqueness, $\exists\eta>0$ such that $x=u$ on $(\tau-\eta,\tau+\eta)\cap[t_0,T_+)$.\\
A crossing would force contact at $e$ without local equality.
\item By monotone boundedness, $\ell\in\R$. If $T_+<+\infty$, local existence at $(T_+,\ell)$ extends $x$, contradiction; hence $T_+=+\infty$. If $F(\ell)\neq0$, then $\exists c>0,t_1$ such that $|F(x(t))|\geq c$ with fixed sign for $t\geq t_1$, so $|x(t)-x(s)|=\left|\int_s^t F(x(r))\,dr\right|\geq c(t-s)$ for $t>s\geq t_1$, contradicting convergence. Thus $F(\ell)=0$.
\end{enumerate}
\hlt{Structural Commentary}: Equilibria block crossings; bounded monotone limits are equilibria.\\
\hlt{Key Insight}: Phase-line signs determine invariant intervals.
\end{proof}
\end{lemma}

\begin{theorem} \label{thm:subcritical-harvest-phase-line} \hlt{Subcritical Harvest Phase Line}\\
Let $h_c\defeq\frac14$, $F_h(x)\defeq(1-x)x-h$, $0\leq h<h_c$, and $x_\pm\defeq\frac12(1\pm\sqrt{1-4h})$.\\
Let $x:[t_0,T_+)\to\R$ be the maximal forward solution of $x'=F_h(x)$, $x(t_0)=x_0$. All limits are as $t\uparrow T_+$.
\begin{enumerate}[label=(\roman*)]
\setlength{\itemsep}{0pt}
\item If $x_0\in\{x_-,x_+\}$, then $x\equiv x_0$.
\item If $x_0<x_-$, then $x'<0$ and $x\to-\infty$.
\item If $x_-<x_0<x_+$, then $x'>0$ and $x\to x_+$.
\item If $x_0>x_+$, then $x'<0$ and $x\to x_+$.
\end{enumerate}
\begin{proof}
Set $I_-\defeq(-\infty,x_-)$, $I_0\defeq(x_-,x_+)$, and $I_+\defeq(x_+,+\infty)$.\\
Since $F_h=-(\,\cdot\,-x_-)(\,\cdot\,-x_+)$, $F_h<0$ on $I_-\cup I_+$ and $F_h>0$ on $I_0$.\\
By Lemma \ref{lem:scalar-phase-line-barrier-limit}(i), $I_-,I_0,I_+$ are invariant.
\begin{enumerate}[label=(\roman*)]
\setlength{\itemsep}{0pt}
\item Since $F_h(x_\pm)=0$, by uniqueness the branch is constant.
\item By invariance, $x([t_0,T_+))\subseteq I_-$, hence $x'<0$.\\
If $x$ had a finite lower limit $\ell$, then by Lemma \ref{lem:scalar-phase-line-barrier-limit}(ii), $\ell\in F_h^{-1}(0)\cap I_-=\emptyset$; hence $x\to-\infty$.
\item By invariance, $x([t_0,T_+))\subseteq I_0$, hence $x'>0$. By Lemma \ref{lem:scalar-phase-line-barrier-limit}(ii), $\ell=\lim x\in[x_0,x_+]\cap F_h^{-1}(0)=\{x_+\}$.
\item By invariance, $x([t_0,T_+))\subseteq I_+$, hence $x'<0$. By Lemma \ref{lem:scalar-phase-line-barrier-limit}(ii), $\ell=\lim x\in[x_+,x_0]\cap F_h^{-1}(0)=\{x_+\}$.
\end{enumerate}
\hlt{Structural Commentary}: Two barriers split $\R$ into three invariant intervals.\\
\hlt{Key Insight}: The upper root attracts from both adjacent sides.
\end{proof}
\end{theorem}

\begin{theorem} \label{thm:critical-harvest-phase-line} \hlt{Critical Harvest Phase Line}\\
Let $h_c\defeq\frac14$, $x_c\defeq\frac12$, $F_h(x)\defeq(1-x)x-h$, and $h=h_c$.\\
Let $x:[t_0,T_+)\to\R$ be the maximal forward solution of $x'=F_h(x)$, $x(t_0)=x_0$. All limits are as $t\uparrow T_+$.
\begin{enumerate}[label=(\roman*)]
\setlength{\itemsep}{0pt}
\item If $x_0=x_c$, then $x\equiv x_c$.
\item If $x_0>x_c$, then $x'<0$ and $x\to x_c$.
\item If $x_0<x_c$, then $x'<0$ and $x\to-\infty$.
\end{enumerate}
\begin{proof}
Since $F_{h_c}=-(\,\cdot\,-x_c)^2$, $F_{h_c}<0$ on $\R\setminus\{x_c\}$ and $F_{h_c}^{-1}(0)=\{x_c\}$.
\begin{enumerate}[label=(\roman*)]
\setlength{\itemsep}{0pt}
\item Since $F_{h_c}(x_c)=0$, by uniqueness the branch is constant.
\item By Lemma \ref{lem:scalar-phase-line-barrier-limit}(i), $x>x_c$, hence $x'<0$. By Lemma \ref{lem:scalar-phase-line-barrier-limit}(ii), $\ell=\lim x\in[x_c,x_0]\cap F_{h_c}^{-1}(0)=\{x_c\}$.
\item By Lemma \ref{lem:scalar-phase-line-barrier-limit}(i), $x<x_c$, hence $x'<0$.\\
If $x$ had a finite lower limit $\ell$, then by Lemma \ref{lem:scalar-phase-line-barrier-limit}(ii), $\ell\in F_{h_c}^{-1}(0)\cap(-\infty,x_c)=\emptyset$; hence $x\to-\infty$.
\end{enumerate}
\hlt{Structural Commentary}: One double barrier remains at $x_c$.\\
\hlt{Key Insight}: The double root attracts from above only.
\end{proof}
\end{theorem}

\begin{theorem} \label{thm:supercritical-harvest-phase-line} \hlt{Supercritical Harvest Phase Line}\\
Let $h_c\defeq\frac14$, $x_c\defeq\frac12$, $F_h(x)\defeq(1-x)x-h$, and $h>h_c$.\\
Let $x:[t_0,T_+)\to\R$ be the maximal forward solution of $x'=F_h(x)$, $x(t_0)=x_0$. Then:
\begin{enumerate}[label=(\roman*)]
\setlength{\itemsep}{0pt}
\item $x'<0$ on $[t_0,T_+)$.
\item $x\to-\infty$ as $t\uparrow T_+$.
\end{enumerate}
\begin{proof}
For $h>h_c$, $F_h=-(\,\cdot\,-x_c)^2-(h-h_c)<0$ on $\R$ and $F_h^{-1}(0)=\emptyset$.
\begin{enumerate}[label=(\roman*)]
\setlength{\itemsep}{0pt}
\item Since $x'=F_h(x)$, $x'<0$.
\item If $x$ had a finite lower limit $\ell$, then by Lemma \ref{lem:scalar-phase-line-barrier-limit}(ii), $\ell\in F_h^{-1}(0)=\emptyset$; hence $x\to-\infty$.
\end{enumerate}
\hlt{Structural Commentary}: No equilibrium barrier remains.\\
\hlt{Key Insight}: Strictly negative velocity precludes a finite limiting zero.
\end{proof}
\end{theorem}

\begin{corollary} \label{cor:critical-logistic-harvest} \hlt{Critical Dimensional Harvest Rate}\\
Let $\kappa,N_{\max}>0$, $H\geq0$, $G_H(N)\defeq\kappa(1-\frac{N}{N_{\max}})N-H$, $H_c\defeq\frac{\kappa N_{\max}}4$. For forward solution $N$, $N(t_0)=N_0$.
\begin{enumerate}[label=(\roman*)]
\setlength{\itemsep}{0pt}
\item $G_H^{-1}(0)\cap(0,+\infty)\neq\emptyset \iff 0\leq H\leq H_c$; $H_c=\max\{H\geq0:G_H^{-1}(0)\cap(0,+\infty)\neq\emptyset\}$.
\item For $0\leq H<H_c$, $G_H^{-1}(0)=\{N_-(H),N_+(H)\}$ with $N_\pm(H)\defeq\frac{N_{\max}}2(1\pm\sqrt{1-\frac{4H}{\kappa N_{\max}}})$; $G_{H_c}^{-1}(0)=\{\frac{N_{\max}}2\}$.
\item If $0\leq H<H_c$ and $N_0>N_-(H)$, then $\lim\limits_{t\uparrow T_+}N(t)=N_+(H)$.
\end{enumerate}
\begin{proof}
Set $x\defeq\frac{N}{N_{\max}}$, $x_0\defeq\frac{N_0}{N_{\max}}$, $\tau\defeq\kappa t$, $h\defeq\frac{H}{\kappa N_{\max}}$, $h_c\defeq\frac14$. Then $\frac{dx}{d\tau}=F_h(x)$, $G_H(N)=\kappa N_{\max}F_h(\frac{N}{N_{\max}})$.
\begin{enumerate}[label=(\roman*)]
\setlength{\itemsep}{0pt}
\item By Proposition \ref{prop:harvested-logistic-equilibria}, $F_h$ has a positive zero iff $0\leq h\leq h_c$.\\
Since $H=\kappa N_{\max}h$, this is $0\leq H\leq H_c$, and the maximal value is $H_c$.
\item By Proposition \ref{prop:harvested-logistic-equilibria}, $F_h^{-1}(0)=\{x_-(h),x_+(h)\}$ for $h<h_c$ and $F_{h_c}^{-1}(0)=\{x_c\}$.\\
After multiplication by $N_{\max}$, $N_\pm(H)=N_{\max}x_\pm(h)$ and $N_{\max}x_c=\frac{N_{\max}}2$.
\item If $N_0>N_-(H)$, then $x_0>\frac{N_-(H)}{N_{\max}}=x_-(h)$.\\
By Theorem \ref{thm:subcritical-harvest-phase-line}(i),(iii),(iv), $\lim\limits_{\tau\uparrow +\infty}x(\tau)=x_+(h)$, hence $\lim\limits_{t\uparrow T_+}N(t)=N_{\max}x_+(h)=N_+(H)$.
\end{enumerate}
\hlt{Structural Commentary}: Scaling sends $H$ to $h=\frac{H}{\kappa N_{\max}}$.\\
\hlt{Key Insight}: $H_c=\frac{\kappa N_{\max}}4$.
\end{proof}
\end{corollary}
